\section{About This Book}\label{app:about}

This companion is a study and prayer guide for keeping Christmastide in the 1962
Roman Rite. It is not an official liturgical book, a critical edition, or a
substitute for your parish calendar, your confessor, or your diocesan ordinance.
Wherever local feasts, transfers, or obligations differ from the universal
pattern described here, follow the calendar and law of your own church.

\subsection{The Liturgical Books and Translations}

The liturgical order followed throughout is the universal Roman Rite of 1962:
the \textit{Missale Romanum} issued by decree of the Sacred Congregation of
Rites on 23~June~1962, governed by Pope St.~John~XXIII's motu proprio
\textit{Rubricarum instructum} (25~July~1960; printed as 23~July~1960 inside the
1962 Missal's own front matter) and the \textit{Codex Rubricarum} promulgated on
26~July~1960 (effective 1~January~1961). Under General Rubric~72~a--b,
\textit{Tempus Natalicium} extends from First Vespers of the Nativity
(24~December evening) through 13~January and is divided into Christmastide
(\textit{Tempus Nativitatis}, First Vespers of 24~December through None of
5~January inclusive) and Epiphanytide (\textit{Tempus Epiphaniae}, First Vespers
of 5~January through 13~January).

For the Latin Collects, Prefaces, and proper \textit{Communicantes}, the text is
the 1962 Roman wording checked against page images of the 1962 typical edition
and collated against the 1862 Ratisbon (Pustet) and 1922 Tours (Mame) Missals.
Because older English hand missals occasionally depart from the 1962 Latin they
print beside them, every English Collect in Appendix~\ref{app:prayers} is drawn
from a verified public-domain translation checked word by word against its 1962
Latin:

\begin{itemize}
\item \textbf{Lasance and Walsh (1945):} Used for Midnight Mass (p.~137), Dawn
Mass (pp.~140--141), St.~Anastasia (p.~141), Day Mass (pp.~143--144), St.~John
(p.~150), the Holy Innocents (p.~154, discarding Lasance's obsolete pre-1960
violet rubric on p.~153), the Sunday within the Octave (pp.~160--161),
St.~Thomas of Canterbury (pp.~157--158), and the Common of Supreme Pontiffs
(p.~1302) for St.~Sylvester~I and St.~Telesphorus (avoiding p.~164, where
Lasance inserts ``Thy'' before ``Supreme Pontiff'').
\item \textbf{Lefebvre (1925):} Used for \textbf{St.~Stephen} (p.~399), because
Lasance (p.~147) omits ``Thy Son'' (\textit{Filium tuum}) from the Collect, and
for the Nativity \textit{Communicantes} (pp.~59--60), which renders both the
night and day variants accurately (neither Lefebvre~1925 nor Lasance~1945
includes St.~Joseph, inserted into the Canon by Pope St.~John~XXIII's decree of
13~November~1962).
\item \textbf{Bute (1908):} Used for the \textbf{Octave Day of the Nativity}
(\textit{Deus, qui salutis aeternae}, p.~315, preserving the complete
Trinitarian conclusion, with the shorter Compline form on p.~211 kept in
\textit{Alma Redemptoris Mater}), because Lasance (pp.~167 and~1403) prints a
misplaced semicolon inside Christ's title (``Our Lord Jesus Christ; Thy Son''),
and for the \textbf{Most Holy Name of Jesus} (p.~365), which preserves the full
Father-addressed conclusion cleanly.
\item \textbf{Further Lasance anomalies avoided:} In the Preface of the Nativity
(p.~798), Lasance mistranslates \textit{cum Thronis et Dominationibus} as ``with
the dominions and powers,'' dropping Thrones; in the Dawn Postcommunion (p.~143),
Lasance attaches ``Whose miraculous birth'' to ``O Lord'' (the Father) rather
than to Christ; and in the St.~John Postcommunion (p.~153), Lasance capitalizes
``Him Whom'' where the Latin refers to St.~John's intercession. None of those
defective renderings is reproduced here.
\end{itemize}

Because a 1961--1962 typical edition of the \textit{Breviarium Romanum} is not
among the books consulted directly here, Latin Office texts---including
\textit{Alma Redemptoris Mater}, \textit{O magnum mysterium}, \textit{Hodie
Christus natus est}, the five \textit{O admirabile commercium} antiphons, the
Christmas and Holy Name hymns, and the Matins patristic lessons---are taken from
the Ratisbon winter Breviary of 1888 (\textit{Breviarium Romanum, pars
hiemalis}, Pustet), paired with John, Marquess of Bute, \textit{The Roman
Breviary} (Winter volume, 1908), and read strictly under the 1960 structural
rubrics (especially the reduction of Sunday within the Octave Matins to one
nocturn of nine psalms and three lessons, RGBR nn.~162~a, 220; the omission of
impeded Scripture lessons, RGBR n.~213; and the three-lesson Saturday Office of
Our Lady, RGBR n.~221~a).

English Scripture quotations follow the Challoner revision of the Douay-Rheims
Bible (using the 1899 American Douay-Rheims text for the Gospel canticles in
Appendix~\ref{app:prayers}), and Latin Scripture follows the Clementine Vulgate.
In the Latin \textit{Benedictus} (Luke~1:70), the electronic \texttt{latVUC}
feed contains a digitization error (\textit{locutum}); the text printed in
Appendix~\ref{app:prayers} restores the received liturgical reading
\textit{sicut locutus est}, verified against Bruno of W\"urzburg (Migne PL~142)
and J.~H.~Blunt (1866). Psalm numbers follow the Vulgate numbering used in the
1962 Roman books, with modern Hebrew numbering in brackets where the two differ.

\paragraph{Historical and textual notes.}
Several scholarly cautions omitted from the narrative chapters for readability
are recorded here: (1) Early historical dates for St.~John Chrysostom's Antioch
homily (c.~386--388), Egeria's Jerusalem witness (c.~385), Duchesne's 25~March
calculation, St.~Francis of Assisi at Greccio (1223), and the Second Council of
Tours (567) rest on the 1908 \textit{Catholic Encyclopedia} articles by
C.~C.~Martindale and J.~Tierney (used in normalized public-domain text). (2) The
quotations from St.~Th\'er\`ese of Lisieux follow T.~N.~Taylor's 1912 English
edition of \textit{Histoire d'une \^ame}, which reflects Mother Agn\`es de
J\'esus's edited text rather than the later unedited critical manuscripts. (3)
William Durandus's exposition of the \textit{Comites Christi} stands in Book~VII,
Chapter~42 of the \textit{Rationale} (Lyon, 1612), while his reflection on
\textit{Dum medium silentium} stands in Book~VI (1568). (4) The attribution of
Midnight Mass and the \textit{Gloria} to Pope St.~Telesphorus rests on the
sixth-century \textit{Liber Pontificalis} tradition echoed by Sicard of Cremona
and the \textit{Roman Martyrology}, not on surviving second-century Roman
liturgical texts. (5) On 1~January, Bute renders \textit{sicut pluvia in vellus}
with ``dew'' rather than literal ``rain,'' and the medieval responsory
\textit{Benedicta et venerabilis es} is the sixth responsory of 1~January rather
than a Christmas Day responsory.

\subsection{The Seven Calendar Shapes of Christmastide}

Because 25~December is fixed on the solar calendar and 25~December through
5~January never crosses 29~February, every year falls into exactly one of seven
recurring shapes determined by the weekday of Christmas Day. The table below was
verified programmatically across 101 consecutive seasons (2020/2021 through
2120/2121; 1,212 liturgical days) and cross-checked with zero discrepancies
against the 1960 \textit{Codex Rubricarum} rules.

{\footnotesize
\begin{longtable}{>{\raggedright\arraybackslash}p{0.10\linewidth}>{\raggedright\arraybackslash}p{0.14\linewidth}>{\raggedright\arraybackslash}p{0.15\linewidth}>{\raggedright\arraybackslash}p{0.12\linewidth}>{\raggedright\arraybackslash}p{0.40\linewidth}}
\toprule
\textbf{25 Dec} & \textbf{Sun in Oct.} & \textbf{Holy Name} & \textbf{Sat BVM} & \textbf{Key 1962 Collisions \& Sample Seasons}\\
\midrule
\endhead
Sunday & None & Mon 2~Jan (fallback) & None & No Sunday within Octave (\textit{Dum medium silentium} omitted); 3~ferias (3--5~Jan); 2022/23, 2033/34, 2039/40.\\
Monday & Sun 31~Dec & Tue 2~Jan (fallback) & None & Sun 24~Dec Vigil replaces Advent~IV (no comm., RG~30~a); Sun 31~Dec omits St.~Sylvester; 2023/24, 2028/29, 2034/35.\\
Tuesday & Sun 30~Dec & Wed 2~Jan (fallback) & Sat 5~Jan & Sat 5~Jan takes \textit{Vultum tuum} (BVM Preface) + St.~Telesphorus comm.\ (Low only); 2029/30, 2035/36, 2040/41.\\
Wednesday & Sun 29~Dec & Sun 5~Jan & Sat 4~Jan & Sun 29~Dec omits St.~Thomas; Sun 5~Jan Holy Name omits St.~Telesphorus even at Low Mass; 2024/25, 2030/31, 2041/42.\\
Thursday & Sun 28~Dec & Sun 4~Jan & Sat 3~Jan & Sun 28~Dec wins over Holy Innocents (comm.\ Low only); Sat 3~Jan \textit{Vultum tuum}; 2025/26, 2031/32, 2036/37.\\
Friday & Sun 27~Dec & Sun 3~Jan & Sat 2~Jan & Sun 27~Dec wins over St.~John (comm.\ Low only); Sat 2~Jan \textit{Vultum tuum}; 2020/21, 2026/27, 2037/38.\\
Saturday & Sun 26~Dec & Sun 2~Jan & None & Sun 26~Dec wins over St.~Stephen (comm.\ Low only); 3~ferias (3--5~Jan); 2021/22, 2027/28, 2032/33.\\
\bottomrule
\end{longtable}
}

Across all 101 seasons (2020/2021--2120/2121), the Sunday within the Octave
occurs 86~times and is absent 15~times; the Most Holy Name falls on a Sunday
58~times and on the 2~January fallback 43~times; Our Lady on Saturday
(\textit{Vultum tuum}) occurs on 57~Saturdays inside 2--5~January; and
Monday--Friday ferias taking the 1~January Mass account for 246~days. Together
with the fixed dates in Table~1, Appendix~\ref{app:daily-course} covers all
1,212~days with no gaps and no overlaps.

\subsection{Canons and Legal Authorities Consulted}

The disciplinary summary in Section~\ref{sec:prayer-home} distinguishes three
legal strata by code, territory, and date (checked as of 27~November~2026):

\begin{itemize}
\item \textbf{Earlier historical discipline:} Council of Trent, Session~XXIV,
\textit{De reformatione matrimonii}, c.~10 (11~November~1563), forbidding solemn
nuptials from Advent~I until the Epiphany; and the Third Plenary Council of
Baltimore's \textit{Manual of Prayers} (New York, 1889, pp.~18--19), recording
nineteenth-century United States practice (Friday abstinence lifted when
Christmas falls on Friday; Sunday Vigil fast anticipated to Saturday; marriages
not solemnized from Advent Sunday until Epiphany).
\item \textbf{1962 universal Latin discipline:} \textit{Codex Iuris Canonici}
(1917), c.~1108 \S\S\,1--3 (solemn nuptial blessing forbidden from Advent~I
through Christmas Day inclusive, lifting on 26~December, unless permitted by the
local Ordinary for just cause), c.~1252 \S\,1 (Friday abstinence, binding on
Fridays inside 26--31~December and 2--5~January), c.~1252 \S\,2 (fast and
abstinence on the Vigil of Christmas), and c.~1252 \S\,4 (cessation on Sundays
and holy days of precept outside Lent, with no vigil anticipation), modified by
the Sacred Congregation of the Council's decree of 3~December~1959 (AAS~51:918)
permitting anticipation of the Vigil fast to 23~December. Because 1917 Code
canons~806 \S\,1 (priests' faculty for three Masses on Christmas Day) and~1247
(the 1917 list of holy days of obligation) were not directly checked in the
legal sources consulted here, those canon numbers are not cited as verified
primary texts.
\item \textbf{Current universal Latin law and United States particular law:}
Pope St.~Paul~VI's apostolic constitution \textit{Paenitemini}
(17~February~1966, Norms~II, IV, VI \S\,1~B), lifting statutory fast and
abstinence from vigils; the \textit{Codex Iuris Canonici} (promulgated
25~January~1983, effective 27~November~1983), cc.~1246--1247 (Sunday Mass and
ten universal holy days of obligation including 25~December and 1~January,
subject to episcopal-conference adaptation under c.~1246 \S\,2) and
cc.~1249--1253 (universal Friday abstinence from age~14 except on solemnities,
leaving Fridays on 26--31~December and 2--5~January penitential because they are
not solemnities); and for the United States, the National Conference of Catholic
Bishops' \textit{Pastoral Statement on Penance and Abstinence}
(18~November~1966, effective 27~November~1966, nos.~17, 22--25), the USCCB
Complementary Norm to cc.~1252--1253 (1983, amended 29~February~1984), and the
2026 USCCB \textit{Liturgical Calendar} miscellaneous notes. Because the
1991/1992 U.S.\ bishops' complementary norm on holy days of obligation
(c.~1246 \S\,2) and 1983 CIC c.~1248 (vigil fulfillment) were not directly
checked in the legal sources consulted here, readers should verify annual
holy-day obligation rules with their diocese or bishops' conference.
\end{itemize}

\subsection{Rights and Provenance}

All quoted Scripture, patristic translations, historical decrees, and received
prayers remain in the public domain or under their own rights as noted in
Appendix~\ref{app:references}; contemporary restricted texts of the Holy See and
the United States Conference of Catholic Bishops are summarized faithfully
without bulk reproduction.
